\chapter*{Abstract}
\addcontentsline{toc}{chapter}{Abstract}

Article~30 of the General Data Protection Regulation (GDPR) requires every controller to maintain a Record of Processing Activities (ROPA),
a tedious task presuming legal expertise small organizations rarely have.
Large Language Model (LLM)-assisted tools have been proposed as a remedy but mostly validated through illustrative use cases rather than controlled studies.
This thesis evaluates ROPAgen, an LLM-powered ROPA tool offering three modes that vary LLM involvement:
Form (checkbox-driven), Ask (user-driven LLM dialogue), and Chat (free-form conversation).
A within-subjects study of $N = 73$ novices used all three modes under Latin Square counterbalancing,
collecting usability, cognitive load, perceived confidence, preferences, and completion time alongside $219$ documents.
Each was scored against an expert-written reference (BLEU, ROUGE, METEOR, BERTScore, SBERT)
and rated by an LLM-as-a-judge on a reference-free rubric covering completeness, faithfulness, legal correctness, and a holistic overall score.
The mode with the least LLM involvement (Form) was the most usable and most preferred.
On document quality, Form led on lexical recall, Chat led on document-level semantic fidelity
and earned the judge's highest holistic verdict, and Ask trailed on every reference-based NLP measure.
Perceived confidence aligned with the judge's holistic verdict but not with the reference-based NLP scores.
On legal correctness the ordering reversed: Chat ranked highest on confidence yet lowest on judge correctness.
The original hypothesis that more LLM involvement yields a better user experience was wrong: less involvement did. More involvement did improve the judge's verdict on the documents, just not the experience the hypothesis was aimed at.
Novices in a domain they do not control need scaffolding, a collaborator, and a safety net, none of which any single mode delivers in full.
The design that follows is a content-typed assignment: enumerable Article~30 sub-records belong to a checkbox interface, scenario-specific ones to conversational extraction.
The LLM is useful as a co-author but not as a teacher; neither a novice nor an LLM alone suffices, and the document that satisfies Article~30 exists only where the two meet.
